% TOP_PROBLEM: {"id": 2744, "title": "Kirby Problem 1.85", "category_id": 7, "category": {"id": 7, "name": "topology", "display_name": "Topology", "description": "Properties preserved under continuous deformations.", "slug": "topology", "order_index": 7, "created_at": "2026-07-31T15:26:25.671Z"}, "status": "open", "problem_number": "KP-1.85", "set_id": 11, "statusQualification": "Uses the character-variety quotient to make the source's conjugacy-class space precise."}
% FIRST_POSTED: 2026-10-01
% ENTRY_KIND: statement_only
% UnsolvedMath ID 2744; source catalogue code KP-1.85
% !TeX program = lualatex
% Definition and sources only; this file is not an LLM solution attempt.
\documentclass[11pt,a4paper]{article}
\usepackage[margin=25mm]{geometry}
\usepackage{fontspec}
\setmainfont{lmroman10-regular.otf}[BoldFont=lmroman10-bold.otf,
 ItalicFont=lmroman10-italic.otf,BoldItalicFont=lmroman10-bolditalic.otf]
\usepackage{amsmath,amssymb,mathtools}
\usepackage{unicode-math}
\setmathfont{latinmodern-math.otf}
\AtBeginDocument{\let\setminus\smallsetminus}
\usepackage{xurl,hyperref}
\hypersetup{colorlinks=true,urlcolor=blue}
\setcounter{secnumdepth}{0}
\setlength{\parindent}{0pt}
\setlength{\parskip}{5pt}
\setlength{\emergencystretch}{3em}
\begin{document}
\section*{Kirby Problem 1.85}\label{2744}
{\small UnsolvedMath ID 2744 · KP-1.85 · Topology · Kirby's Problems in Low-Dimensional Topology}
\subsection{Definitions and mathematical statement}
Let $K\subset S^3$ be a smooth knot whose complement admits a complete finite-volume hyperbolic metric. Set $G=\pi_1(S^3\setminus K)$ and let $X(G)$ be its $\operatorname{PSL}_2(\mathbb C)$ character variety: the algebraic quotient of the representation variety by conjugation. Equivalently, use characters of semisimple representations, so nonclosed conjugacy orbits do not create distinct points of this quotient.

The hyperbolic holonomy gives a discrete faithful representation $\rho_{\mathrm{hyp}}:G\to\operatorname{PSL}_2(\mathbb C)$. Fix either orientation convention for the hyperbolic structure, and denote the irreducible algebraic component containing its character by $X_0(K)$, the canonical component. Regard $\operatorname{SO}(3)\cong\operatorname{PSU}(2)$ as the compact subgroup of $\operatorname{PSL}_2(\mathbb C)$.

Does $X_0(K)$ contain an arc of characters represented by homomorphisms into $\operatorname{SO}(3)$? Here an arc means the image of a continuous injection $[0,1]\to X_0(K)$ in the usual complex topology, all of whose points have such compact-group representatives. In particular, these must be genuinely different characters, not a family obtained merely by conjugating one representation.

The question is universal over hyperbolic knots. The representations on the arc need not themselves be discrete or faithful. Their characters must lie on the specified component, rather than on an arbitrary component of the knot character variety; this is the essential connection with hyperbolic holonomy.
\subsection{Short English statement}
Must the canonical character component of every hyperbolic knot contain a genuine arc of rotation-group characters?
\subsection{Sources}
\begin{itemize}
\item[S1] ULAM.AI and the UnsolvedMath Contributors, supplied problems.json, record 2744 (KP-1.85), Kirby's Problems in Low-Dimensional Topology. Catalogue identity and source transcription; CC BY 4.0.
\url{https://huggingface.co/datasets/ulamai/UnsolvedMath}.
\item[S2] K3: A New Problem List in Low-Dimensional Topology, Problem 1.85. Expanded formulation of the supplied catalogue statement.
\url{https://math.berkeley.edu/sites/default/files/surv-295-ruberman-watermarked-author-pdf.pdf}.
\end{itemize}
\end{document}
